% Generated draft: requires mathematical and source review.
\documentclass[11pt]{article}
\usepackage[margin=25mm]{geometry}
\usepackage{fontspec}
\setmainfont{Latin Modern Roman}
\usepackage{amsmath,amssymb,mathtools,unicode-math}
\setmathfont{Latin Modern Math}
\usepackage{enumitem,needspace,xurl,hyperref}
\setlength{\parindent}{0pt}
\setlength{\parskip}{5pt}
\setlength{\emergencystretch}{3em}
\newcommand{\N}{\mathbb{N}}
\newcommand{\Z}{\mathbb{Z}}
\newcommand{\Q}{\mathbb{Q}}
\newcommand{\R}{\mathbb{R}}
\newcommand{\C}{\mathbb{C}}
\newcommand{\F}{\mathbb{F}}
\newcommand{\A}{\mathbb{A}}
\newcommand{\PP}{\mathbb{P}}
\newcommand{\E}{\mathbb{E}}
\newcommand{\eps}{\varepsilon}
\newcommand{\dd}{\,\mathrm{d}}
\begin{document}
\section*{UnsolvedMath ID: 30003194}
\textbf{Source problem code:} catalogue entry
\subsection{Definitions and mathematical statement}
Let \(k\) be a finite field, let \(\overline{k}\) be an algebraic closure, let \(X\) be a variety over \(k\), and let \(\overline{x}\) be a geometric base point of \(X_{\overline{k}}\). For a prime \(\ell\), put
\[
G_\ell=\pi_1^{\mathrm{\acute et},\ell}(X_{\overline{k}},\overline{x}),
\]
where the superscript denotes the maximal pro-\(\ell\) quotient of the geometric étale fundamental group. Write
\[
\Lambda_\ell=\mathbb Z_\ell[[G_\ell]]
=\varprojlim_{m,U}(\mathbb Z/\ell^m\mathbb Z)[G_\ell/U],
\qquad
I=\ker\bigl(\Lambda_\ell\xrightarrow{\mathrm{aug}}\mathbb Z_\ell\bigr),
\]
where \(m\geq 1\), \(U\) ranges over open normal subgroups of \(G_\ell\), and augmentation sends every group element to \(1\). Assume that every augmentation quotient \(I^n/I^{n+1}\) is torsion-free as a \(\mathbb Z_\ell\)-module, for \(n\geq 0\).

The supplied question asks whether there exists a real number \(r>0\) such that the convergent group algebra denoted in the source record by
\[
B_r=\mathbb Q_\ell[[G_\ell]]^{\leq\ell^{-r}}
\]
admits a set \(E\) of Frobenius eigenvectors whose linear span is dense in \(B_r\). With the intended scalar field and topology understood, the requested conclusion is
\[
\overline{\operatorname{span}(E)}^{\,B_r}=B_r.
\]
The question is posed for arbitrary \(X\), \(k\), and prime \(\ell\) subject to the stated augmentation hypothesis; no uniform choice of \(r\) is requested.

\textit{Unresolved qualifications.} The supplied record does not define the convergence condition represented by \({}^{\leq\ell^{-r}}\), its topology, or the precise rational completed algebra used to construct \(B_r\). Accordingly, \(B_r\) and the displayed closure retain their source notation rather than specifying a complete analytic construction. The record also does not specify a Frobenius operator on this based algebra, the convention for Frobenius, or whether eigenvectors are taken over \(\mathbb Q_\ell\) or after scalar extension. A Frobenius eigenvector would mean a nonzero vector \(v\) satisfying \(F(v)=\lambda v\) for the intended operator \(F\) and an eigenvalue in the intended field; these choices remain unresolved. No restriction \(\ell\neq\operatorname{char}(k)\), geometric connectedness assumption, or base-point compatibility with Frobenius is supplied, so none is imposed here. The original extraction refers to a completed enveloping algebra, whereas the corrected statement specifies the convergent group algebra above; their intended relationship is not supplied. This is therefore a statement-only reconstruction with unresolved formulation details.
\subsection{Short English statement}
For a variety over a finite field and a prime \(\ell\), assume that all augmentation quotients of the completed integral group algebra of its geometric pro-\(\ell\) fundamental group are torsion-free. The question asks whether some positive convergence parameter gives a rational convergent group algebra with a dense linear span of Frobenius eigenvectors. The convergence topology, Frobenius action, and eigenvalue field require clarification. The supplied record labels the problem open; that status is not independently verified here.
\subsection{Sources}
\begin{itemize}
\item Local UnsolvedMath registry: supplied complete source record, permanent ID \(30003194\), exact source problem code \texttt{catalogue entry}, title ``Dense Frobenius Eigenvectors in Completed Enveloping Algebras,'' in the Oberwolfach Reports Open Problems collection. No individual proposer is identified; the proposed year is recorded as 2016.
\item Supplied report citation: \emph{Arithmetic Geometry (2017)}, Oberwolfach Reports, DOI 10.4171/owr/2016/38, pp.\ 2171--2224. Source URL and supplied literature source, labelled ``OWR Arithmetic Geometry report'': \url{https://doi.org/10.4171/owr/2016/38}. The record describes these reports as documenting mathematical workshops and their open-problem discussions. The citation's 2017 designation and the record's 2016 workshop/proposal designation are retained without resolving their relationship.
\item Supplied statement-verification metadata labels the statement \texttt{corrected\_verified} and says that the official report was checked to restore the maximal pro-\(\ell\) group, coefficient rings, augmentation hypothesis, convergent group algebra, and dense-eigenvector conclusion. This is a supplied verification claim, not an independent source check performed for this draft.
\item Literature assessment dated exactly \texttt{2026-08-21}: the supplied curation labels the problem open and describes the extraction as an incomplete conjectural extension of Theorem 5, with partial cases in the cited report and no general resolution reported there. This is supplied dated curation, not independently verified. The content and hypotheses of Theorem 5 and the partial cases are not included in the record; the supplied curation requests independent verification of time-sensitive status and cited claims.
\end{itemize}
\end{document}
